\documentclass[10pt,a4paper]{amsart}
\usepackage{amsmath,amssymb,amsthm,graphicx}
\usepackage{verbatim}
\newenvironment{reference}{\comment}{\endcomment}
%\newenvironment{reference}{}{}
\newenvironment{slogan}{\comment}{\endcomment}
\newenvironment{history}{\comment}{\endcomment}

% For commutative diagrams we use Xy-pic
\usepackage[all]{xy}

% We use 2cell for 2-commutative diagrams.
\xyoption{2cell}
% Theorem environments.
%
\theoremstyle{plain}
\newtheorem{theorem}[subsection]{Theorem}
\newtheorem{proposition}[subsection]{Proposition}
\newtheorem{lemma}[subsection]{Lemma}

\theoremstyle{definition}
\newtheorem{definition}[subsection]{Definition}
\newtheorem{example}[subsection]{Example}
\newtheorem{exercise}[subsection]{Exercise}
\newtheorem{situation}[subsection]{Situation}

\theoremstyle{remark}
\newtheorem{remark}[subsection]{Remark}
\newtheorem{remarks}[subsection]{Remarks}

\numberwithin{equation}{subsection}

% Macros
%
\def\lim{\mathop{\mathrm{lim}}\nolimits}
\def\colim{\mathop{\mathrm{colim}}\nolimits}
\def\Spec{\mathop{\mathrm{Spec}}}
\def\Hom{\mathop{\mathrm{Hom}}\nolimits}
\def\Ext{\mathop{\mathrm{Ext}}\nolimits}
\def\SheafHom{\mathop{\mathcal{H}\!\mathit{om}}\nolimits}
\def\SheafExt{\mathop{\mathcal{E}\!\mathit{xt}}\nolimits}
\def\Sch{\mathit{Sch}}
\def\Mor{\mathop{\mathrm{Mor}}\nolimits}
\def\Ob{\mathop{\mathrm{Ob}}\nolimits}
\def\Sh{\mathop{\mathit{Sh}}\nolimits}
\def\NL{\mathop{N\!L}\nolimits}
\def\CH{\mathop{\mathrm{CH}}\nolimits}
\def\proetale{{pro\text{-}\acute{e}tale}}
\def\etale{{\acute{e}tale}}
\def\QCoh{\mathit{QCoh}}
\def\Ker{\mathop{\mathrm{Ker}}}
\def\Im{\mathop{\mathrm{Im}}}
\def\Coker{\mathop{\mathrm{Coker}}}
\def\Coim{\mathop{\mathrm{Coim}}}

% Boxtimes
%
\DeclareMathSymbol{\boxtimes}{\mathbin}{AMSa}{"02}

%
% Macros for moduli stacks/spaces
%
\def\QCohstack{\mathcal{QC}\!\mathit{oh}}
\def\Cohstack{\mathcal{C}\!\mathit{oh}}
\def\Spacesstack{\mathcal{S}\!\mathit{paces}}
\def\Quotfunctor{\mathrm{Quot}}
\def\Hilbfunctor{\mathrm{Hilb}}
\def\Curvesstack{\mathcal{C}\!\mathit{urves}}
\def\Polarizedstack{\mathcal{P}\!\mathit{olarized}}
\def\Complexesstack{\mathcal{C}\!\mathit{omplexes}}
% \Pic is the operator that assigns to X its picard group, usage \Pic(X)
% \Picardstack_{X/B} denotes the Picard stack of X over B
% \Picardfunctor_{X/B} denotes the Picard functor of X over B
\def\Pic{\mathop{\mathrm{Pic}}\nolimits}
\def\Picardstack{\mathcal{P}\!\mathit{ic}}
\def\Picardfunctor{\mathrm{Pic}}
\def\Deformationcategory{\mathcal{D}\!\mathit{ef}}
\usepackage{cleveref}
\crefname{originalreference}{original reference}{original references}
\Crefname{originalreference}{Original reference}{Original references}
\makeatletter
\newlabel{situation-setup}{{02QL}{source}{}{sourcepointer-1}{}}
\newlabel{situation-setup@cref}{{[originalreference][0][]02QL}{[1][0][]source}{}{}{}}
\makeatother
\crefname{originallemma}{original lemma}{original lemmas}
\Crefname{originallemma}{Original lemma}{Original lemmas}
\makeatletter
\newlabel{lemma-compute-koszul}{{0FEE}{source}{}{sourcepointer-2}{}}
\newlabel{lemma-compute-koszul@cref}{{[originallemma][0][]0FEE}{[1][0][]source}{}{}{}}
\makeatother
\crefname{originallemma}{original lemma}{original lemmas}
\Crefname{originallemma}{Original lemma}{Original lemmas}
\makeatletter
\newlabel{lemma-chern-classes-dual}{{0FA5}{source}{}{sourcepointer-3}{}}
\newlabel{lemma-chern-classes-dual@cref}{{[originallemma][0][]0FA5}{[1][0][]source}{}{}{}}
\makeatother
\crefname{originallemma}{original lemma}{original lemmas}
\Crefname{originallemma}{Original lemma}{Original lemmas}
\makeatletter
\newlabel{lemma-vectorbundle}{{02TY}{source}{}{sourcepointer-4}{}}
\newlabel{lemma-vectorbundle@cref}{{[originallemma][0][]02TY}{[1][0][]source}{}{}{}}
\makeatother
\crefname{originallemma}{original lemma}{original lemmas}
\Crefname{originallemma}{Original lemma}{Original lemmas}
\makeatletter
\newlabel{lemma-restrict-to-open}{{02RX}{source}{}{sourcepointer-5}{}}
\newlabel{lemma-restrict-to-open@cref}{{[originallemma][0][]02RX}{[1][0][]source}{}{}{}}
\makeatother
\crefname{originallemma}{original lemma}{original lemmas}
\Crefname{originallemma}{Original lemma}{Original lemmas}
\makeatletter
\newlabel{lemma-top-chern-class}{{0FA9}{source}{}{sourcepointer-6}{}}
\newlabel{lemma-top-chern-class@cref}{{[originallemma][0][]0FA9}{[1][0][]source}{}{}{}}
\makeatother
\title{The top localized Chern class of the section of a projective bundle vanishes below the rank and is a factorial multiple of the vector bundle pullback}
\author{The Stacks Project Authors}
\date{stacks-project:chow.tex:local-snapshot-20261012; GNU FDL 1.2 or later}
\begin{document}
\maketitle
\noindent\textit{Editorial note.} All delivered mathematical text is the author's own, copied from the pinned Stacks Project chapter file chow.tex; no statement or proof step was rewritten, reordered or removed.\par\medskip
\noindent\textit{Editorial note.} Preamble provenance: the pinned chapter file opens with an input of the shared preamble and carries no preamble text of its own, so the project's single Stacks-wide preamble staged at sources/stacks-chow/preamble.tex is excerpted verbatim as two spans, lines 14-18 (Xy-pic setup) and lines 154-217 (the native theorem declarations and mathematical macros). Nothing was added, removed or redefined; the wrapper class and the amsmath/amssymb/amsthm/graphicx base packages are supplied by the assembler.\par\medskip
\noindent\textit{Editorial note.} The counted statement and proof contain no cross reference, no citation and no included file, so the delivered body is self contained.\par\medskip
\noindent\textit{Editorial note.} Canonical location of the counted proposition: Stacks tag 0FEF (sources/stacks-chow/pinned-original-source.tex line 12542).\par\medskip
\noindent\textit{Allowed standard prerequisites (editorial).}\begin{itemize}
\item The Koszul resolution of the section of a projective bundle and the computation of its Chern classes. Same chapter, Lemma labelled lemma-compute-koszul, cited at proof lines 12577-12580. Shows i\_* O\_X is perfect, gives the vanishing below r and computes c\_r as a factorial multiple.
\item Chern classes of a dual and of a bundle with a regular section. Same chapter, Lemma labelled lemma-chern-classes-dual, cited at proof lines 12582-12584, and Lemma labelled lemma-top-chern-class, cited at proof lines 12604-12606. Rewrites the class into the form whose zero locus is i(X) and evaluates it on integral subschemes.
\item Cycle classes on a vector bundle and the restriction sequence to its complement. Same chapter, Lemma labelled lemma-vectorbundle, cited at proof lines 12592-12594, and Lemma labelled lemma-restrict-to-open, cited at proof lines 12594-12596. Decomposes alpha as pi\textasciicircum{}* beta + gamma with gamma supported off C.
\end{itemize}\medskip
\ifdefined\hypertarget\hypertarget{sourcepointer-1}{}\fi
\noindent\textit{Source reference.} Original reference 02QL: sources/stacks-chow/pinned-original-source.tex, source lines 1372--1372. Same-chapter pointer into the author environment situation at source line 1372; invoked at proof line None of the counted proof. It is retained only as a navigation pointer; every load-bearing step is inside the delivered span.\par\medskip
\ifdefined\hypertarget\hypertarget{sourcepointer-2}{}\fi
\noindent\textit{Source reference.} Original lemma 0FEE: sources/stacks-chow/pinned-original-source.tex, source lines 12500--12500. Same-chapter pointer into the author environment lemma at source line 12500; invoked at proof line 12574 of the counted proof. It is retained only as a navigation pointer; every load-bearing step is inside the delivered span.\par\medskip
\ifdefined\hypertarget\hypertarget{sourcepointer-3}{}\fi
\noindent\textit{Source reference.} Original lemma 0FA5: sources/stacks-chow/pinned-original-source.tex, source lines 8520--8520. Same-chapter pointer into the author environment lemma at source line 8520; invoked at proof line 12580 of the counted proof. It is retained only as a navigation pointer; every load-bearing step is inside the delivered span.\par\medskip
\ifdefined\hypertarget\hypertarget{sourcepointer-4}{}\fi
\noindent\textit{Source reference.} Original lemma 02TY: sources/stacks-chow/pinned-original-source.tex, source lines 7062--7062. Same-chapter pointer into the author environment lemma at source line 7062; invoked at proof line 12593 of the counted proof. It is retained only as a navigation pointer; every load-bearing step is inside the delivered span.\par\medskip
\ifdefined\hypertarget\hypertarget{sourcepointer-5}{}\fi
\noindent\textit{Source reference.} Original lemma 02RX: sources/stacks-chow/pinned-original-source.tex, source lines 2949--2949. Same-chapter pointer into the author environment lemma at source line 2949; invoked at proof line 12594 of the counted proof. It is retained only as a navigation pointer; every load-bearing step is inside the delivered span.\par\medskip
\ifdefined\hypertarget\hypertarget{sourcepointer-6}{}\fi
\noindent\textit{Source reference.} Original lemma 0FA9: sources/stacks-chow/pinned-original-source.tex, source lines 8750--8750. Same-chapter pointer into the author environment lemma at source line 8750; invoked at proof line 12605 of the counted proof. It is retained only as a navigation pointer; every load-bearing step is inside the delivered span.\par\medskip
\begin{lemma}
\label{lemma-compute-section}
Let $(S, \delta)$ be as in Situation \ref{situation-setup}. Let $X$
be a scheme locally of finite type over $S$. Let $\mathcal{C}$
be a locally free $\mathcal{O}_X$-module of rank $r$. Consider the
morphisms
$$
X = \underline{\text{Proj}}_X(\mathcal{O}_X[T])
\xrightarrow{i}
E = \underline{\text{Proj}}_X(\text{Sym}^*(\mathcal{C})[T])
\xrightarrow{\pi}
X
$$
Then $c_t(i_*\mathcal{O}_X) = 0$ for $t = 1, \ldots, r - 1$ and in
$A^0(C \to E)$ we have
$$
p^* \circ \pi_* \circ c_r(i_*\mathcal{O}_X) = (-1)^{r - 1}(r - 1)! j^*
$$
where
$j : C \to E$ and $p : C \to X$ are the inclusion and structure
morphism of the vector bundle
$C = \underline{\Spec}(\text{Sym}^*(\mathcal{C}))$.
\end{lemma}

\begin{proof}
The canonical map $\pi^*\mathcal{C} \to \mathcal{O}_E(1)$ vanishes
exactly along $i(X)$. Hence the Koszul complex on the map
$$
\pi^*\mathcal{C} \otimes \mathcal{O}_E(-1) \to \mathcal{O}_E
$$
is a resolution of $i_*\mathcal{O}_X$. In particular we see that
$i_*\mathcal{O}_X$ is a perfect object of $D(\mathcal{O}_E)$
whose Chern classes are defined. The vanishing of $c_t(i_*\mathcal{O}_X)$
for $t = 1, \ldots, t - 1$ follows from Lemma \ref{lemma-compute-koszul}.
This lemma also gives
$$
c_r(i_*\mathcal{O}_X) = - (r - 1)!
c_r(\pi^*\mathcal{C} \otimes \mathcal{O}_E(-1))
$$
On the other hand, by Lemma \ref{lemma-chern-classes-dual} we have
$$
c_r(\pi^*\mathcal{C} \otimes \mathcal{O}_E(-1)) =
(-1)^r c_r(\pi^*\mathcal{C}^\vee \otimes \mathcal{O}_E(1))
$$
and $\pi^*\mathcal{C}^\vee \otimes \mathcal{O}_E(1)$ has a section $s$
vanishing exactly along $i(X)$.

\medskip\noindent
After replacing $X$ by a scheme locally of finite type over $X$,
it suffices to prove that both sides of the equality have the
same effect on an element $\alpha \in \CH_*(E)$. Since $C \to X$
is a vector bundle, every cycle class on $C$ is of the form $p^*\beta$
for some $\beta \in \CH_*(X)$ (Lemma \ref{lemma-vectorbundle}).
Hence by Lemma \ref{lemma-restrict-to-open}
we can write $\alpha = \pi^*\beta + \gamma$ where $\gamma$
is supported on $E \setminus C$. Using the equalities above
it suffices to show that
$$
p^*(\pi_*(c_r(\pi^*\mathcal{C}^\vee \otimes \mathcal{O}_E(1)) \cap [W])) =
j^*[W]
$$
when $W \subset E$ is an integral closed subscheme which
is either (a) disjoint from $C$ or (b) is of the form $W = \pi^{-1}Y$
for some integral closed subscheme $Y \subset X$.
Using the section $s$ and Lemma \ref{lemma-top-chern-class} we find
in case (a) $c_r(\pi^*\mathcal{C}^\vee \otimes \mathcal{O}_E(1)) \cap [W] = 0$
and in case (b)
$c_r(\pi^*\mathcal{C}^\vee \otimes \mathcal{O}_E(1)) \cap [W] = [i(Y)]$.
The result follows easily from this; details omitted.
\end{proof}

\end{document}
